% Generated draft: requires mathematical and source review.
\documentclass[11pt]{article}
\usepackage[margin=25mm]{geometry}
\usepackage{fontspec}
\setmainfont{Latin Modern Roman}
\usepackage{amsmath,amssymb,mathtools,unicode-math}
\setmathfont{Latin Modern Math}
\usepackage{enumitem,needspace,xurl,hyperref}
\setlength{\parindent}{0pt}
\setlength{\parskip}{5pt}
\setlength{\emergencystretch}{3em}
\newcommand{\N}{\mathbb{N}}
\newcommand{\Z}{\mathbb{Z}}
\newcommand{\Q}{\mathbb{Q}}
\newcommand{\R}{\mathbb{R}}
\newcommand{\C}{\mathbb{C}}
\newcommand{\F}{\mathbb{F}}
\newcommand{\A}{\mathbb{A}}
\newcommand{\PP}{\mathbb{P}}
\newcommand{\E}{\mathbb{E}}
\newcommand{\eps}{\varepsilon}
\newcommand{\dd}{\,\mathrm{d}}
\begin{document}
\section*{UnsolvedMath ID: 800798}
\textbf{Source problem code:} catalogue entry
\subsection{Definitions and mathematical statement}
An affine hyperplane in \(\R^4\) has the form
\[
H_i=\{x\in\R^4:\langle a_i,x\rangle=b_i\},\qquad a_i\in\R^4\setminus\{0\},\quad b_i\in\R.
\]
For four such hyperplanes and a sign vector \(\sigma\in\{-1,1\}^4\), define the corresponding open region by
\[
C_\sigma=\{x\in\R^4:\sigma_i(\langle a_i,x\rangle-b_i)>0\text{ for every }i=1,\ldots,4\}.
\]
These regions are pairwise disjoint and cover \(\R^4\setminus\bigcup_{i=1}^4H_i\).

For the discrete formulation below, use the convention that a finite set is in general position if every subset of at most five points is affinely independent. This convention is explicit here because the supplied record does not define general position. Interpret separation as requiring every point to lie outside the four hyperplanes.

Under these conventions, the discrete question is: for every integer \(n\geq1\) and every set \(P\subset\R^4\) of \(16n\) points in general position, do there exist \(a_1,\ldots,a_4\in\R^4\setminus\{0\}\) and \(b_1,\ldots,b_4\in\R\) such that
\[
P\cap\bigcup_{i=1}^4H_i=\varnothing,
\qquad
\lvert P\cap C_\sigma\rvert=n
\quad\text{for every }\sigma\in\{-1,1\}^4?
\]
Thus the sixteen subsets must arise from the sixteen sign regions of the same four hyperplanes.

The source also asks whether every ``object'' in \(\R^4\) can be partitioned into sixteen ``equal parts'' by four hyperplanes. The supplied record specifies neither the class of objects nor the meaning of equality. It also does not prescribe how portions lying on a cutting hyperplane are handled. Consequently, this broader question cannot be given a unique precise mathematical statement from the supplied evidence; no measure, regularity assumption, or notion of congruence is imposed here.
\subsection{Short English statement}
The discrete question asks whether four hyperplanes can place any \(16n\) points in general position in \(\R^4\) into sixteen sign regions containing exactly \(n\) points each. The draft explicitly uses affine general position and requires the cutting hyperplanes to avoid the points. The broader question about dividing every object into equal parts remains underspecified.
\subsection{Sources}
\begin{itemize}
\item Local UnsolvedMath registry, permanent ID \(800798\), exact source code \texttt{catalogue entry}, title ``Branko Grunbaum in Many Dimensions --- Question 5,'' in the AMR Open Problem Lists collection. The record attributes the problem to Kahle and gives the proposed year as 2019.
\item Supplied source list: Kahle, \emph{Branko Grunbaum in Many Dimensions: Open Problems} (2019), Question 5, \url{https://arxiv.org/abs/1901.08622}. The record reports access on 2026-10-16 and extraction by \texttt{arxiv\_latex}; these are supplied provenance details, not an independent verification of the source.
\item The local record labels the problem ``open.'' This is a supplied status label, not an independently verified mathematical or literature assessment.
\end{itemize}
\end{document}
